\section{Introduction}\label{sec:introduction}

Let \(X\) be a compact irreducible holomorphic symplectic (IHS) Kähler
manifold of dimension \(2n\). A nontrivial fibration of \(X\) is a proper
surjective holomorphic map \(f:X\to B\) with connected fibers onto a normal
complex space, with \(0<\dim B<2n\). In the projective setting, Matsushita
proved that \(\dim B=n\) and that a smooth general fiber is an abelian
variety \cite{Matsushita1999,Matsushita2001}. His equidimensionality theorem
shows that every irreducible component of every fiber is Lagrangian once
the general fiber is Lagrangian \cite[Theorem~1]{Matsushita2000}.
For compact Kähler IHS total spaces, Greb--Lehn's Proposition~2.2 shows
that the base is projective of dimension \(n\) and that every fiber
component is Lagrangian \cite{GrebLehn2014}. The projective-space base conjecture asks
whether this normal base must be isomorphic to \(\PP^n\).

Hwang proved the projective-space conclusion when \(X\) is projective and
\(B\) is smooth \cite[Theorem~1.2]{Hwang2008}. Greb and Lehn extended the
smooth-base theorem to compact Kähler IHS total spaces
\cite[Theorem~1.1]{GrebLehn2014}. These results reduce the base problem to
the possible singularities of \(B\).

The conclusion is also known for the \(K3^{[n]}\), generalized Kummer,
\(\mathrm{OG6}\), and \(\mathrm{OG10}\) deformation types without assuming
the base smooth. A common nef formulation of the earlier isotropic-divisor
and rational-fibration results is recorded in
\cite[Theorem~1.3]{DebarreHuybrechtsMacriVoisin2024}: a nef line bundle \(L\)
whose first Chern class is primitive and isotropic for the
Beauville--Bogomolov--Fujiki form on such a compact IHS manifold
admits a holomorphic Lagrangian fibration \(g:X\to\PP^n\) with
\(L\simeq g^*\cO_{\PP^n}(1)\). For an existing fibration \(f:X\to B\), choose
an ample line bundle \(H\) on \(B\) and write \(f^*H\simeq L^r\) with
\(c_1(L)\) primitive. This class remains nef and isotropic, so the theorem
applies to \(L\); connected fibers identify
the graded section rings of \(H\), \(L^r\), and \(\cO_{\PP^n}(r)\).
Taking their \(\operatorname{Proj}\) identifies \(B\) with \(\PP^n\).

In dimension four, Ou reduced the normal surface base of a projective IHS
fibration to \(\PP^2\) or the specified Fano surface
\(S^{\mathrm n}(E_8)\), with one \(E_8\) singularity and two nodal rational
anticanonical curves \cite[Theorem~1.2]{Ou2019}. Huybrechts and Chenyang Xu
excluded the second case \cite[Theorem~0.1]{HuybrechtsXu2022}.
Müller and Zheng Xu subsequently proved a local theorem: a projective
Lagrangian germ from a smooth holomorphic symplectic manifold has smooth
base if the normal analytic base has only isolated quotient singularities
\cite[Theorem~1.3]{MullerXu2026}. Their Corollary~1.4 treats such germs in
total dimension four without a singularity assumption, and their
Theorem~1.5 treats terminal threefold bases in total dimension six.
Their Example~4.1 exhibits a nonisolated quotient stratum for a noncompact
symplectic total space; excluding these strata in the compact setting
therefore requires additional global geometry.

For projective IHS total spaces, the cohomological constraint is already
that of projective space: Shen--Yin compute the rational intersection
cohomology of \(B\) as one-dimensional in degrees \(0,2,\ldots,2n\) and
zero in odd degrees \cite[Theorem~0.4(a)]{ShenYin2022}. This does not itself establish
smoothness. Our proof first establishes smoothness and then applies Hwang's
theorem.

\begin{theorem}\label{thm:main}
Let \(n\geq1\), and let \(X\) be a compact irreducible holomorphic symplectic
Kähler manifold of dimension \(2n\), with symplectic form \(\sigma\).
Suppose that
\[
                       f:X\longrightarrow B
\]
is a projective surjective holomorphic morphism with connected fibers,
where \(B\) is a normal projective variety of dimension \(n\).
Assume that every irreducible component of every fiber has dimension
\(n\) and that the restriction of \(\sigma\) to its smooth locus is zero.
Then \(B\) is isomorphic to \(\PP^n_{\C}\).
\end{theorem}

Here projectivity of \(f\) means the existence of a relatively ample
holomorphic line bundle. Together with projectivity of \(B\), it makes the
original \(X\) projective. Compact Kähler deformations, which can be
nonprojective, enter later in the proof. The theorem imposes no smoothness
condition on \(B\), no additional lower bound on \(b_2(X)\), and no
restriction on the deformation type. It gives a positive answer to the
projective-space base conjecture for the stated projective fibrations.

The proof uses the companion paper \cite{StrongSYZ2026}.
Its Theorem~1.1 supplies semiampleness of a nonzero
integral nef isotropic line bundle on the given compact Kähler IHS manifold
and the associated connected Lagrangian fibration; we restate it as
Theorem~\ref{thm:strong-syz}. Its Lemma~4.2 supplies a polarized family for
a prescribed primitive ample class, with polarization a positive multiple
of that class. It also supplies a torsion-free finite-index subgroup of the
integral lattice isometries fixing the class and preserving the chosen
period-domain component. This subgroup acts freely, the quotient is smooth
and quasi-projective, and a local marked lift of the family's algebraic
period map is submersive at some point. The family monodromy lies in that subgroup.
Lemma~\ref{secsys:actual-monodromy} of this article proves that its image
has finite index there. Section~\ref{sec:second-system}
uses these inputs to construct a second Lagrangian fibration.

Once \(B\simeq\PP^n\) has been established, Kim--Oguiso--Shinder's theorem
excludes codimension-one multiple fibers
\cite[Theorems~1.1 and~3.4]{KimOguisoShinder2026}. Kamenova--Verbitsky also
prove this conclusion over \(\PP^n\)
\cite[Theorems~1.5 and~6.7]{KamenovaVerbitsky2026}. Thus, for every prime
divisor \(D\subset B\), writing \(f^*D=\sum_i a_iE_i\) with distinct prime
divisors \(E_i\), one has \(\gcd_i a_i=1\). Equivalently, there is no
equality \(f^*D=mE\) with an integral effective divisor \(E\) and \(m>1\).
This is a consequence of the base identification.

\subsection{Proof strategy}

The proof has two stages. First we rule out germs that are not analytic
quotients of smooth germs by finite groups. We then eliminate the
stabilizers of the remaining quotient germs. The second stage uses both
the original fibration and a finite cover by an abelian variety constructed
from deformations that retain \(B\). This cover is available without
assuming that \(B\) has quotient singularities.

Suppose \(B\) has a nonquotient germ. Sections~\ref{sec:local-tensor}
and~\ref{sec:stratum} use finite local covers and reflexive cotangent tensors
to find a nonempty closed normal projective stratum \(S_0\subset B\) of
dimension less than \(n\). A smooth Deligne--Mumford stack \(\cS\) with coarse
space \(S_0\) records the finite stabilizers on its local charts.
Sections~\ref{sec:base-change} and~\ref{sec:euler} prove the obstruction
\[
                         \chi_{\mathrm{top}}(S_0)=0,
\]
where this is the ordinary Euler characteristic of the coarse space.
The numerical argument compares Euler characteristics after pullback to
\(X\) with alternating cotangent contributions on the local charts.
This comparison uses ordinary coefficient base change in a suitable
reduction to positive characteristic; Frobenius growth and Riemann--Roch
force the relevant values to be zero. The local tensor classification
then allows coherent coefficient modules to isolate the trivial character
of the full stabilizer. This recovers the ordinary Euler characteristic
of \(S_0\), including when some stabilizers act trivially on the stratum.

The tensor recognition in this construction uses Deligne's criterion
\cite[Theorem~7.1]{Deligne1990}. The differential identity behind the
base-change formula uses the strictness and decomposition of projective
direct images of Hodge modules \cite{Saito1988,Saito2016}. These tools enter the
Euler obstruction; the later exclusions use the global geometry of the
original fibration and its deformations.

Choose a very ample line bundle \(H\) on \(B\) and put \(e=c_1(f^*H)\).
Matsushita's deformation theorem gives a relatively generated line bundle
and a family of Lagrangian maps throughout the nearby locus where \(e\)
remains of type \((1,1)\)
\cite[Theorem~1.2 and Corollary~1.3]{MatsushitaDeformations2016}.
The targets in that result may vary. Section~\ref{sec:deformation} proves
that their normal Stein targets can all be identified with the original
\(B\), obtaining a map to \(B\times D\) over that entire nearby locus.
Verbitsky's degenerate-twistor construction already keeps a Lagrangian
fibration holomorphic over its original base along a complex line
\cite[Theorem~1.10]{Verbitsky2015}. Bogomolov--Déev--Verbitsky develop
this construction for closed base forms of type \((2,0)+(1,1)\)
\cite[Proposition~3.2 and Theorem~3.3]{BogomolovDeevVerbitsky2022}.
Soldatenkov--Verbitsky give a formulation for normal analytic bases and
show, for compact IHS fibrations, that these deformations are Zariski
locally trivial over the projective base and arise from twists
\cite[Section~2.2, Theorem~3.1 and Corollary~3.3]{SoldatenkovVerbitsky2025}.
The simultaneous fixed target over the full nearby \(e\)-Hodge locus is
the additional assertion needed here.

Lemma~\ref{lem:betti-bound} gives \(b_2(X)\geq4\). If \(b_2(X)\geq5\), period
variation in the fixed-base family forces the Hodge cohomology of \(\cS\)
to be diagonal. Consequently
\[
             \chi_{\mathrm{top}}(S_0)=\sum_p h^{p,p}(\cS)>0,
\]
contradicting the obstruction. When \(b_2(X)=4\), we distinguish the
isotrivial family of polarized smooth fibers from a family with nonzero
variation. Section~\ref{sec:isotrivial} treats the isotrivial case: its
ramification and reflection arguments reduce a possible stratum to an
elliptic curve, where a connection on a negative line bundle gives the
contradiction. In the other case, Section~\ref{sec:second-system} constructs
a projective deformation \(f':X'\to B\) with a second, varying Lagrangian
fibration \(g:X'\to C\). Section~\ref{sec:variation} uses its monodromy and
the symmetric period cubic to force the same diagonal Hodge conclusion.
The cubic is the Donagi--Markman condition on the smooth polarized locus
\cite[Section~7.2, Lemmas~7.1--7.2 and Remark~7.3]{DonagiMarkman1996};
Deligne's fixed-part theorem supplies the monodromy invariants used with it
\cite[Theorem~4.1.1(ii) and Corollary~4.1.2]{Deligne1971}.
These three cases rule out every nonquotient germ.

Starting from the original fibration and its fixed-base family, the
second-system construction also supplies a cover in every allowed rank.
On a smooth fiber \(A\) of \(g\), the mixed Fujiki formula makes the nef
line bundle \(f'^*H|_A\) have positive top self-intersection, hence it is
ample on the abelian variety \(A\). Therefore \(f'|_A\) is finite, and
equality of dimensions makes it surjective. This is a finite cover by an
abelian variety.
After the first stage has proved that \(B\) has quotient singularities,
Section~\ref{sec:quotient} applies this cover and the original fibration.
On the canonical smooth stack over \(B\), the direct-image identities
produce a single translation vector on \(A\) whose projection to the directions
fixed by each nonidentity stabilizer element never vanishes on its fixed locus.
This use of equivariant higher direct images has a local
precedent in Müller--Xu's combination of those identities, cohomology and
base change for the flat fibration, and equivariant Lefschetz vanishing
\cite[Theorem~3.1]{MullerXu2026}. Here the finite relation associated with
the cover records pairs of points of \(A\) together with an isomorphism
between their stack images. A component containing a nonidentity
automorphism over a pair \((a,a)\) has a proper difference image in \(A\).
Choose such a component for which that image has maximal dimension.
Polarization and
ramification produce another such component with a larger difference
image, giving a contradiction. Thus \(B\) is smooth, and Hwang's theorem
completes the proof.

Figure~\ref{fig:proof} displays the three exclusions in the first stage
and the cover available in every allowed rank.

\begin{figure}[htbp]
\centering
\begin{tikzpicture}[
  box/.style={draw=MidnightBlue!70,rounded corners=2pt,
    fill=MidnightBlue!3,align=center,inner sep=6pt,font=\small},
  line/.style={semithick,draw=MidnightBlue!80},
  arr/.style={line,-{Stealth[length=2mm]}}]
\node[box,text width=11.6cm] (zero) at (0,0)
  {A nonquotient germ gives a closed stratum \(S_0\) with
   \(\chi_{\mathrm{top}}(S_0)=0\)\\
   (Sections~\ref{sec:local-tensor}--\ref{sec:euler})};
\node[box,text width=3.65cm,minimum height=2cm] (large) at (-4.5,-2.3)
  {\(b_2\geq5\)\\Fixed-base deformation\\and period variation\\
   (Section~\ref{sec:deformation})};
\node[box,text width=3.65cm,minimum height=2cm] (iso) at (0,-2.3)
  {\(b_2=4\), isotrivial\\Ramification and an\\elliptic connection\\
   (Section~\ref{sec:isotrivial})};
\node[box,text width=3.65cm,minimum height=2cm] (vary) at (4.5,-2.3)
  {\(b_2=4\), varying\\Fixed-base deformation\\with a second fibration\\and symmetric cubic\\
   (Sections~\ref{sec:second-system}--\ref{sec:variation})};
\node[box,text width=5.7cm,minimum height=1.45cm] (quot) at (-3.4,-5)
  {Each case excludes the stratum:\\
   \(B\) has only quotient singularities};
\node[box,text width=5.7cm,minimum height=1.45cm] (cover) at (3.4,-5)
  {Second-system construction:\\a finite cover of \(B\) by an abelian variety\\
   (Section~\ref{sec:second-system}, all \(b_2\geq4\))};
\node[box,text width=11.6cm] (smooth) at (0,-7)
  {Eliminate quotient stabilizers (Section~\ref{sec:quotient})\\
   \(B\) is smooth; Hwang's theorem gives \(B\simeq\PP^n\)};
\draw[arr] (zero.south west)--(large.north);
\draw[arr] (zero.south)--(iso.north);
\draw[arr] (zero.south east)--(vary.north);
\draw[line] (large.south)--(-4.5,-3.75)--(0,-3.75);
\draw[line] (iso.south)--(0,-3.75);
\draw[line] (vary.south)--(4.5,-3.75)--(0,-3.75);
\draw[arr] (0,-3.75)--(quot.north);
\draw[arr] (quot.south)--(smooth.north west);
\draw[arr] (cover.south)--(smooth.north east);
\end{tikzpicture}
\caption{The first stage rules out a nonquotient germ in three cases.
The second uses both quotient singularities and a finite cover by an
abelian variety.}
\label{fig:proof}
\end{figure}
